\section{Preliminary interface evidence}
This section reports bounded observations from the existing extraction, not final corpus-wide results. The current 21-record working set combines 18 retained candidates from the earlier S1--S3 audit with three subsequent S4--S5 studies. It includes methodological and complementary evidence, and is not a set of 21 demonstrations of uncertainty continuity.

\subsection{Building attributes to material quantities}
The semisupervised material-intensity study provides probabilistic representations that can support downstream stock modelling \citep{U03}. This capability is distinct from demonstrating that classification uncertainty and quantity uncertainty are jointly propagated through a complete inventory. The feature-based stock estimation audit identifies a further distinction between reported predictive error and uncertainty retained during aggregation of point estimates \citep{U05}. These observations motivate separate extraction of parameter uncertainty, prediction uncertainty and the aggregation rule.

\subsection{Material stocks to temporal flows}
The lifespan study examines the sensitivity of stock and demolition accounts to lifetime assumptions, while the probabilistic dynamic MFA study explicitly connects uncertain inputs with stock and flow outputs \citep{U11,U12}. This interface changes the temporal allocation of material quantities. Its uncertainty cannot be characterized solely by the accuracy of the total stock: flow timing and spatial aggregation can alter the relevance of lifetime assumptions. These studies provide useful mechanisms for the synthesis, without establishing that every upstream observation error has been retained.

\subsection{Recovery and environmental consequences}
The multicycle renovation framework links reuse feasibility and recovery-related inputs with circularity and environmental assessment \citep{U19}. The extraction distinguishes shared inputs and compatible scenarios from a fully specified joint stochastic transfer between the branches. The circularity score is not itself an input passed into LCA. The reported Monte Carlo calculations do not document all distributional dependencies or paired draws needed to reproduce a complete cross-branch uncertainty chain.

Damage-to-environmental-impact studies provide complementary mappings between physical states, interventions and carbon consequences \citep{U20,U21}. They cannot automatically be treated as evidence for uncertain material recovery yields. Within the current working set, the S3--S4 connection remains insufficiently established for a claim of continuous uncertainty propagation from released quantities to recoverable quantities and quality. This is a bounded finding, not evidence that the wider literature contains no such mechanism.

\subsection{Domain-specific DT evidence}
The broader S1--S5 DT retrieval and the narrower uncertainty-gated DT retrieval are separate corpora. Their screening does not yet support a final architecture-frequency analysis or a validated domain-wide gap map. This subsection will be developed from interface-level full-text extraction rather than from keyword occurrence or retrieval counts. Candidate architectures will be assessed for entity linkage, model coupling and actual update behavior before their capabilities are attributed to a functioning twin.
